Se generaron dos conjuntos de datos: (1) ordenamiento y (2) multiplicación de matrices, usando scripts en Python.

\subsection*{1. Ordenamiento}
Arreglos de enteros:  
- Tamaños $N$: $10^1$, $10^3$, $10^5$, $10^7$.  
- Tipos: aleatorio, ascendente, descendente.  
- Dominios: $D1=\{0,\ldots,9\}$, $D7=\{0,\ldots,10^7\}$.  
- Réplicas: 3 (a,b,c) por configuración.  

\subsection*{2. Multiplicación de Matrices}
Matrices cuadradas $N\times N$: $16$, $64$, $256$, $1024$.  
- Tipos: densas, dispersas ($\sim90\%$ ceros), diagonales.  
- Dominios: $D0=\{0,1\}$, $D10=\{0,\ldots,9\}$.  
- Réplicas: 3 pares (a,b,c) por configuración.  

\medskip
Esta variedad de matrices permite analizar si la estructura interna de los datos (por ejemplo, la dispersión) tiene un impacto en el rendimiento práctico de los algoritmos, algo que la complejidad teórica no considera.
\medskip